% -----------------------------------------------------------
\section{1873}
% -----------------------------------------------------------

	% ----------------------
	\subsection{Brigham Young}
	% ----------------------
		\label{EMS-1873-BY}

		% ----------------------
		\subsubsection{Introduction to 1873}
		% ----------------------
		
			sdf
			
		% -----
		\subsubsection{Discussion, 9 June 1873}
		% -----
			\label{EMS-1873-BY-9-JUN-1873}
			% \label{EMS-18XX-BY-DAY-3LETTERMONTH-YEAR}
			
			\begin{description}
					\item[Print Sources]
				\end{description}

					\begin{tabular}{p{0.5in}p{1in}p{0.5in}p{1in}}
					CDBY	& ????		& JD 		& ???? \\
					\end{tabular}
				
				\begin{description}
					\item[Online Access] %\href{https://archive.org/details/diaryofljohnnutt01nutt/page/n135/mode/2up}{Here}
					% \item[Analysis] \hyperref[XX]{Here}
				\end{description}
				
				\begin{description}
					\item[Background]
				\end{description}
				
					\phantom{.}
				
				\begin{description}
					\item[Original Source]
				\end{description}
				
					\textit{Salt Lake School of the Prophets 1867--1883}, edited by Devery S. Anderson, Signature Books, 2018, 377--378
					% brief description of original source material, with reference to JSP or somewhere else for more info
					% Background material: it might be helpful to have a note that says whether a given document was presented as a speech, was written in a journal, etc.
					% Also a comment here about how reliable the source might be, or what shape it is in, if there are large omissions or parts that are hard to understand, and so on.
					% Also a comment about when I transcribed it, whether I've checked it more than once, and so on
				
				\begin{description}
					\item[Text]
				\end{description}
				
					The School of the Prophets met in the Old Tabernacle at 4 o'clock p.m. Present Prest Brigham Young, Presiding. Prayer by William Hickenlooper. Roll called. 

					Prest D[aniel] H. Wells remarked that Prest Young touched on certain principles and doctrine yesterday, and particularly that doctrine pertaining to Adam being our Father \& our God, and thought this was a proper place to call for the minds of the brethren and learn whether we as a school endorse the doctrine.
					
					He bore a powerful testimony to the truth of the doctrine, remarking that if ever he had received a testimony of any doctrine in this Church he had of the truth of this. The Endowments plainly teach it and the Bible \& other revelations are full\sout{y} of it.
					
					The \sout{principle} doctrine was approved or endorsed by Henry Grow, D[imick] B. Huntingdon, \& Joseph F. Smith. The latter read a portion of a revelation given to the Church, \^{}page 2[01] Doc. \& Cov.\^{} affirming that Michael or Adam is the Father of all---the Prince of all, and stated that the enunciation of that doctrine, gave him great joy. A[lonzo] H. Raleigh said he had never heard any one dispute the doctrine, to him it was perfectly natural, as was every other principle of the Gospel when understood.
					
					Prest Young queried whether the brethren thought he was too liberal in launching out on this doctrine before the Gentiles. He was positive on the truth of the doctrine, but thought we should be cautious about preaching on doctrines unless we fully understand them by the power of the Spirit, then they commend themselves to the hearts of our hearers. Spoke of the vain theories of men with regard to the Great first Cause. Said there were many revelations given to him that he did not receive from the Prophet Joseph [Smith]. He did not receive them through the Urim and Thummim as Joseph did but when he did receive them he knew of their truth as much as it was possible for \sout{them} him to do of any truth. He thought it advisable to adjourn the School, as it was a busy time and only about one third of the brethren present.
					
					John Lyon fully endorsed the \sout{princip} doctrine, but asked explanation why the Scriptures seemed to put Jesus Christ on an equal footing with the Father.
					
					Prest Young replied that the writers of those scriptures wrote according to their best language and understanding.
					
					Geo B. Wallace endorsed the doctrine---\& narrated a dream he received in answer to prayer, shewing that many worlds were made to God because he had many wives.
					
					Joseph F. Smith said \sout{he did not believe} ``a man could \sout{enter} \^{}not attain\^{} unto \sout{the} \^{}a\^{} fullness of \^{}the\^{} glory \^{}of the Father\^{}, who did not take to himself more than one wife.''